\begin{thebibliography}{99}
\bibitem{Shamai1990} Shamai S. Capacity of a pulse amplitude modulated direct detection photon channel. IEE Proc I. 1990;137: 424--430. doi: 10.1049/ip-i-2.1990.0056
\bibitem{Nikitin2009} Nikitin AP, Stocks NG, Morse RP, McDonnell MD. Neural population coding is optimized by discrete tuning curves. Phys Rev Lett. 2009;103: 138101. doi: 10.1103/physrevlett.103.138101
\bibitem{DBS2021} Dytso A, Barletta L, Shamai S. Properties of the support of the capacity-achieving distribution of the amplitude-constrained Poisson noise channel. IEEE Trans Inf Theory. 2021;67: 7050--7066. doi: 10.1109/tit.2021.3111836
\bibitem{BDS2024} Barletta L, Dytso A, Shamai S. Improved bounds on the number of support points of the capacity-achieving input for amplitude constrained Poisson channels. In: Proceedings of the IEEE International Symposium on Information Theory (ISIT); 2024. pp. 3630--3635. doi: 10.1109/isit57864.2024.10619622
\bibitem{Cao2014} Cao J, Hranilovic S, Chen J. Capacity-achieving distributions for the discrete-time Poisson channel---Part I: General properties and numerical techniques. IEEE Trans Commun. 2014;62: 194--202. doi: 10.1109/tcomm.2013.112513.130142
\bibitem{Paper1} Firdaus IMTH. How many atoms, where, and when: Predicting the discrete optima of peak-power capacity, minimax estimation and nonparametric maximum likelihood under Gaussian noise. Companion paper; 2026. Available from: \url{https://github.com/Tristan-tech-ai/how-many-atoms}
\bibitem{KostalShinomoto2016} Kostal L, Shinomoto S. Efficient information transfer by Poisson neurons. Math Biosci Eng. 2016;13: 509--520. doi: 10.3934/mbe.2016004
\bibitem{Favano2026} Favano A, Baniasadi M, Zieder I, Barletta L, Dytso A. The binomial channel: on capacity, optimal inputs, and beta-binomial approximation. arXiv:2607.02683 [Preprint]. 2026.
\bibitem{Barlow1961} Barlow HB. Possible principles underlying the transformations of sensory messages. In: Rosenblith WA, editor. Sensory communication. Cambridge (MA): MIT Press; 2012. pp. 216--234. doi: 10.7551/mitpress/9780262518420.003.0013%
\bibitem{Laughlin1981} Laughlin SB. A simple coding procedure enhances a neuron's information capacity. Z Naturforsch C. 1981;36: 910--912. doi: 10.1515/znc-1981-9-1040
\bibitem{McDonnellStocks2008} McDonnell MD, Stocks NG. Maximally informative stimuli and tuning curves for sigmoidal rate-coding neurons and populations. Phys Rev Lett. 2008;101: 058103. doi: 10.1103/physrevlett.101.058103
\bibitem{Wen2009} Wen B, Wang GI, Dean I, Delgutte B. Dynamic range adaptation to sound level statistics in the auditory nerve. J Neurosci. 2009;29: 13797--13808. doi: 10.1523/jneurosci.5610-08.2009
\bibitem{Escabi2003} Escabí MA, Miller LM, Read HL, Schreiner CE. Naturalistic auditory contrast improves spectrotemporal coding in the cat inferior colliculus. J Neurosci. 2003;23: 11489--11504. doi: 10.1523/jneurosci.23-37-11489.2003
\bibitem{BerryMeister1998} Berry MJ 2nd, Meister M. Refractoriness and neural precision. J Neurosci. 1998;18: 2200--2211. doi: 10.1523/jneurosci.18-06-02200.1998
\bibitem{Heeringa2024} Heeringa AN. Single-unit data for sensory neuroscience: Responses from the auditory nerve of young-adult and aging gerbils. Sci Data. 2024;11: 411. doi: 10.1038/s41597-024-03259-3
\bibitem{Holm1979} Holm S. A simple sequentially rejective multiple test procedure. Scand J Stat. 1979;6: 65--70.
\bibitem{Shmueli2005} Shmueli G, Minka TP, Kadane JB, Borle S, Boatwright P. A useful distribution for fitting discrete data: revival of the Conway-Maxwell-Poisson distribution. J R Stat Soc Ser C Appl Stat. 2005;54: 127--142. doi: 10.1111/j.1467-9876.2005.00474.x
\bibitem{Watson1922} Watson GN. A treatise on the theory of Bessel functions. Cambridge: Cambridge University Press; 1922.
\end{thebibliography}
